% ============================================================================
%  Precautionary Saving in Europe --- Project Documentation
%  Compile with:  pdflatex documentation.tex   (run twice for references)
%  Font: Palatino (mathpazo, with matching math).
% ============================================================================
\documentclass[11pt,a4paper]{article}

% ---- Fonts: Palatino for text and math -------------------------------------
\usepackage[T1]{fontenc}
\usepackage[utf8]{inputenc}
\usepackage{mathpazo}            % Palatino text + matching maths (Hermann Zapf)
\linespread{1.05}                % Palatino reads better with a touch more leading
\usepackage{microtype}

% ---- Layout ----------------------------------------------------------------
\usepackage[a4paper,margin=2.6cm]{geometry}
\usepackage{booktabs}
\usepackage{array}
\usepackage{amsmath,amssymb}
\usepackage{graphicx}
\graphicspath{{figures/}{./}}
\usepackage{caption}
\captionsetup{font=small,labelfont=bf}
\usepackage{enumitem}
\setlist{itemsep=2pt,topsep=3pt}
\usepackage{xcolor}
\definecolor{codebg}{gray}{0.96}
\definecolor{rule}{gray}{0.5}

% ---- Code / paths listing style --------------------------------------------
\usepackage{listings}
\lstset{
  basicstyle=\ttfamily\small,
  backgroundcolor=\color{codebg},
  frame=single, rulecolor=\color{rule},
  breaklines=true, columns=fullflexible,
  showstringspaces=false, keepspaces=true,
  xleftmargin=4pt, xrightmargin=4pt, framesep=4pt,
}

\usepackage[hidelinks,colorlinks=true,linkcolor=black!70!blue,
            urlcolor=black!60!blue,citecolor=black]{hyperref}

% ---- Convenience macros ----------------------------------------------------
\newcommand{\code}[1]{\texttt{#1}}
\newcommand{\dataset}[1]{\textsc{#1}}
% Robust figure inclusion: shows a placeholder if the PNG is absent.
\newcommand{\projfig}[3]{% {width-fraction}{file-without-ext}{caption}
  \begin{figure}[htbp]\centering
    \IfFileExists{figures/#2.png}
      {\includegraphics[width=#1\linewidth]{#2}}
      {\fbox{\parbox{0.78\linewidth}{\centering\itshape figure not found ---
        run the pipeline (\code{02\_make\_figures.py}) to generate it.}}}
    \caption{#3}\label{fig:#2}
  \end{figure}}
% Figures produced by the follow-the-money extension.
\newcommand{\extfig}[3]{% {width-fraction}{file-without-ext}{caption}
  \begin{figure}[htbp]\centering
    \IfFileExists{extension_follow_money/figures/#2.png}
      {\includegraphics[width=#1\linewidth]{extension_follow_money/figures/#2}}
      {\fbox{\parbox{0.78\linewidth}{\centering\itshape figure not found ---
        run \code{extension\_follow\_money/run.sh} to generate it.}}}
    \caption{#3}\label{fig:#2}
  \end{figure}}

\title{\textbf{Precautionary Saving in Europe}\\[2pt]
       \large Data, Pipeline, and Econometric Methodology}
\author{Project documentation}
\date{\today}

\begin{document}
\maketitle
\thispagestyle{empty}

\begin{abstract}
\noindent
This document explains, step by step, the empirical project behind the
hypothesis that the euro area's elevated household saving rate is
\emph{substantially precautionary} --- a response to geopolitical and economic
uncertainty rather than to income or wealth alone. It documents (i) the
reproducible three-stage software pipeline, (ii) every data series used, with an
explicit justification of its source, frequency, length, units, and treatment,
and (iii) the econometric strategy in full detail: stationarity classification,
the ARDL bounds test for a long-run relationship, Granger causality, and a
structural VAR with orthogonalised impulse responses. A final section turns to
the \emph{composition} of saving --- a ``follow-the-money'' test that, unlike the
level analysis, can separate a yield-chasing reaction to ECB rates from a
precautionary one, and finds the post-2022 reallocation is far more likely the
former. Reported figures come from a clean end-to-end run; they are produced live
and may be revised by the data providers, so they should be regenerated before
citation.
\end{abstract}

\tableofcontents
\bigskip
\hrule
\bigskip

% ============================================================================
\section{Overview and research question}
% ============================================================================

\paragraph{Hypothesis.} The euro-area household saving rate has stayed well above
its pre-pandemic norm even after pandemic forced-saving and the income/wealth
shocks unwound. We test whether this persistence is \emph{precautionary}: driven
by elevated uncertainty, it should
\begin{enumerate}[label=(\Alph*)]
  \item move \emph{with} uncertainty over time (the time-series leg);
  \item have risen \emph{more} where the Russia/energy shock bit hardest (the
        cross-country leg);
  \item be concentrated among households who can \emph{afford} to save, while the
        fear itself weighs most on those who cannot (the distributional leg).
\end{enumerate}
Each leg maps to a chart and, for leg~(A), to a battery of time-series tests.
The guiding methodological principle is that \emph{co-movement is not
causation}: uncertainty, policy rates, and inflation all rose together after
2022, so the descriptive charts can only show association. Section~\ref{sec:econ}
is the attempt to isolate the precautionary channel formally.

\paragraph{What the documentation covers.} Section~\ref{sec:pipeline} describes
the reproducible pipeline. Section~\ref{sec:data} justifies every data choice.
Section~\ref{sec:descriptive} walks through the descriptive charts.
Section~\ref{sec:econ} is the detailed econometric methodology and results.
Sections~\ref{sec:interpretation}--\ref{sec:limits} synthesise and qualify the
findings. The appendix lists the file tree, exact commands, and a data
dictionary.

% ============================================================================
\section{The reproducible pipeline}\label{sec:pipeline}
% ============================================================================

The project is deliberately split into three scripts by \emph{function}, so that
each stage can be run, audited, and re-run independently. Data collection (slow,
network-bound, non-deterministic across vintages) is separated from figure
rendering (fast, deterministic, offline) and from econometric analysis
(deterministic, offline).

\begin{center}
\renewcommand{\arraystretch}{1.25}
\begin{tabular}{@{}p{4.4cm} p{4.0cm} p{6.2cm}@{}}
\toprule
\textbf{Stage / script} & \textbf{Reads} & \textbf{Writes} \\
\midrule
\code{01\_collect\_data.py} & live APIs (Eurostat, FRED, GPR, OECD) &
  \code{data/*.csv} (tidy series) \\
\code{02\_make\_figures.py} & \code{data/*.csv} &
  \code{figures/*.png} (charts A--F) \\
\code{03\_econometrics.py} & \code{data/*.csv} &
  \code{data/econometrics\_results.md}, \code{figures/A2*, A3*} \\
\bottomrule
\end{tabular}
\end{center}

\paragraph{One-command reproduction.} The wrapper \code{run\_all.sh} installs the
pinned dependencies and runs the three stages in order:
\begin{lstlisting}
bash run_all.sh        # pip install, then run 01 -> 02 -> 03 (classic Python)
\end{lstlisting}

\paragraph{Classic Python, no virtual environment.} The project runs directly on
a standard system Python ($\ge 3.11$); there is nothing to create or activate.
\code{run\_all.sh} installs the requirements into the chosen interpreter
(\code{python3} by default; override with the \code{PYTHON} variable) and then
runs the three stages. Dependencies are pinned in \code{requirements.txt} so the
numerical results are stable across machines. If the default \code{python3} is a
Homebrew/``externally-managed'' build that refuses a global \code{pip install},
point \code{PYTHON} at a python.org interpreter or add \code{--user} to the pip
line.

\paragraph{Network robustness.} Downloads use the \code{requests} library, which
ships its own certificate bundle (\code{certifi}); this avoids the
\code{SSL: CERTIFICATE\_VERIFY\_FAILED} errors that the standard-library URL
reader raises on a stock macOS Python. Each data pull is wrapped in
\code{try/except}: a failed source prints a clear message and leaves any existing
CSV from a previous run untouched, so a transient outage in one provider never
corrupts the rest of the dataset. The Caldara--Iacoviello index is distributed
as a legacy \code{.xls} file, which is why \code{xlrd} is a pinned dependency.

% ============================================================================
\section{Data}\label{sec:data}
% ============================================================================

This section is the heart of the ``why this data'' question. For each series we
state the \emph{source} and exact dataset code, the \emph{frequency} and why it
is appropriate, the \emph{length/coverage}, the \emph{units}, and any
\emph{transformation} and \emph{caveat}. All series are free and require no API
key. Spans below are from a representative run and will extend as providers
release new vintages.

% ----------------------------------------------------------------------------
\subsection{The target: euro-area household saving rate (quarterly)}
% ----------------------------------------------------------------------------
\begin{description}[leftmargin=1.4em,style=nextline]
\item[Source / code] Eurostat sector accounts, \dataset{nasq\_10\_ki} (key
  indicators), item \code{SRG\_S14\_S15}, sector \code{S14\_S15} (households and
  NPISH), seasonally and calendar adjusted (\code{SCA}), geography \code{EA20}.
\item[Frequency] Quarterly. This is the natural and \emph{binding} frequency: the
  saving rate is only published quarterly, and it is the dependent variable, so
  all higher-frequency regressors are aggregated down to it (Sec.~\ref{sec:freq})
  rather than the reverse.
\item[Coverage] 1999Q1--2025Q4 ($\approx$108 quarters); median level $\approx
  13.2\%$, latest $\approx 14.4\%$.
\item[Units] Gross household saving as a percentage of gross disposable income
  (the gross saving rate). We use the \emph{gross} rate because it is the
  headline Eurostat definition with the longest, most comparable series; this is
  noted whenever cross-source comparisons arise.
\item[Why this exact series] The \dataset{nasq\_10\_ki} table contains many
  ratios (saving rate, investment rate, profit share, $\ldots$) under different
  \code{na\_item} codes. The collector \emph{never averages across indicators};
  it auto-selects the single dimension combination whose values look like a
  household saving rate---preferring a saving-named \code{na\_item} and a level
  inside a plausible band of $5\%$--$30\%$, closest to a typical $\sim$13\%---and
  refuses any candidate outside that band. This guards against a silent
  mis-selection that, in an earlier version, blended unrelated ratios into a
  nonsensical $\sim$35\% ``saving rate''. The datasets \dataset{nasa\_10\_ki} and
  \dataset{teina500} are tried as fallbacks.
\end{description}

% ----------------------------------------------------------------------------
\subsection{Uncertainty proxies}
% ----------------------------------------------------------------------------
Precautionary saving theory is about \emph{perceived risk}, which is
unobservable; we proxy it two ways and let the data adjudicate
(Sec.~\ref{sec:timing}).
\begin{description}[leftmargin=1.4em,style=nextline]
\item[Geopolitical Risk (GPR)] Caldara \& Iacoviello, a newspaper-based index of
  geopolitical risk. Monthly, here from 1999-01 onward (the index itself extends
  back to 1900). It is the preferred proxy for this project because the framing
  is explicitly geopolitical (it spikes cleanly at the February 2022 invasion).
  Distributed as a keyless \code{.xls}.
\item[EU Economic Policy Uncertainty (EPU)] FRED series \code{EUEPUINDXM}
  (Baker--Bloom--Davis methodology for Europe). Monthly, 1999-01 onward. A
  broader ``economic-policy'' notion of uncertainty, used as the robustness
  alternative to GPR.
\end{description}
\emph{Frequency rationale.} Both are monthly; uncertainty can move sharply
within a quarter, so a monthly source preserves the timing of spikes, which we
then average to quarterly to match the target.

% ----------------------------------------------------------------------------
\subsection{Confounders (controls)}
% ----------------------------------------------------------------------------
Two variables co-move with uncertainty and independently affect saving; omitting
them would bias any estimate of the precautionary channel.
\begin{description}[leftmargin=1.4em,style=nextline]
\item[ECB policy rate] FRED \code{ECBDFR} (deposit facility rate); fallbacks
  \code{ECBMRRFR}, then a 3-month interbank rate. Published \emph{daily}
  ($\approx$10{,}000 observations from 1999); we take the quarterly mean. A higher
  policy rate raises the reward to saving, a competing explanation for the
  post-2022 rise.
\item[HICP inflation] Eurostat \dataset{prc\_hicp\_manr}, all-items \code{CP00},
  \code{EA20}, year-on-year \%. Monthly, from 2000-12. High inflation both erodes
  real incomes and can itself raise precaution; it must be held constant.
\end{description}

% ----------------------------------------------------------------------------
\subsection{Cross-country cross-section (leg B)}
% ----------------------------------------------------------------------------
\begin{description}[leftmargin=1.4em,style=nextline]
\item[Saving by country] Eurostat \dataset{tec00131} (item \code{SRG\_S14\_S15}),
  \emph{annual}, 21 EU countries, 2014--2025. Annual is the only frequency at
  which a comparable cross-country household saving rate exists; for a
  cross-section that is sufficient.
\item[Energy shock] Eurostat \dataset{prc\_hicp\_manr} restricted to energy
  (\code{NRG}), monthly YoY \%. For each country we take the \emph{peak} monthly
  energy inflation during calendar 2022, a single scalar measuring how hard the
  energy shock hit.
\item[Derived join] \code{scatter\_energy\_vs\_saving.csv} pairs each country's
  energy peak with its \emph{change} in saving from 2019 to the 2023/24 average.
  Using the change (not the level) nets out structural North--South differences
  in saving habits, the more defensible test of the shock's effect.
\end{description}

% ----------------------------------------------------------------------------
\subsection{Distribution (leg C)}
% ----------------------------------------------------------------------------
There is no single official, annual, euro-area saving-rate-by-quintile series, so
the project triangulates.
\begin{description}[leftmargin=1.4em,style=nextline]
\item[Eurostat ICW] \dataset{icw\_sr\_03} (experimental ``Income, Consumption and
  Wealth'' statistics), median saving rate by income quintile (\code{QU1}--%
  \code{QU5}), unit \code{PC\_DI} (\% of disposable income), geography \code{EA}.
  Compiled only at multi-year reference points ($\sim$2010/2015/2020), so it is
  read as \emph{structural} (the shape of the distribution), not as an annual
  path.
\item[OECD EG DNA] Expert Group on Distributional National Accounts, attempted
  via the OECD SDMX API for an annual by-quintile panel. \textbf{Status:} the
  queried flow returns data without a usable quintile dimension, so chart~E is
  not produced. The pull is retained, clearly flagged, for when the correct flow
  id is confirmed; it is a known data-availability limitation, not a code error.
\item[ECB CES expectations] Expected unemployment 12 months ahead by income
  group, from published ECB Consumer Expectations Survey figures. These are not
  in any bulk API, so they are entered as documented constants in
  \code{01\_collect\_data.py} and written to
  \code{ces\_unemp\_expectations.csv}; refresh from the latest CES release before
  publishing.
\end{description}

% ----------------------------------------------------------------------------
\subsection{Frequency alignment}\label{sec:freq}
% ----------------------------------------------------------------------------
The econometrics needs one common frequency. Because the dependent variable (the
saving rate) is quarterly, \emph{quarterly is the binding frequency}: every
monthly or daily regressor is collapsed to a quarterly mean
($\text{resample}(\text{QS}).\text{mean}()$), and series are inner-joined on
date. Aggregating down (averaging) rather than interpolating up avoids inventing
within-quarter variation that the data do not contain. After alignment the
estimation sample is \textbf{2000Q4--2025Q4 (101 quarters)}; the start is set by
the shortest series (HICP inflation, from 2000-12) and the end by the
target. The quarterly mean also lightly smooths the monthly uncertainty proxies,
which is innocuous for the relationships we study.

% ============================================================================
\section{Descriptive evidence}\label{sec:descriptive}
% ============================================================================

The charts are intentionally simple; they motivate the formal tests but do not
substitute for them.

\projfig{0.92}{A_saving_vs_uncertainty}{\textbf{Leg A.} Top: the euro-area
household saving rate against its 2012--19 norm, with the post-2022 excess
shaded and the February 2022 invasion marked. Bottom: the candidate
drivers---uncertainty, the ECB rate, and HICP inflation---standardised to
$z$-scores. They visibly rise together, which is precisely why eyeballing cannot
isolate the precautionary channel.}

\projfig{0.78}{B_energy_vs_saving_scatter}{\textbf{Leg B.} Each country's peak
2022 energy inflation (the shock) against its change in saving 2019$\to$2023/24.
Colour marks editorial Russia/energy exposure; the dashed line is the OLS fit.
The slope is positive but noisy.}

\projfig{0.62}{B2_saving_by_country_bar}{\textbf{Leg B, levels.} Household saving
rate by country in the latest year. Germany and the Nordics sit at the top---but
they were structurally high savers before the war, which is the confound that
motivates using the \emph{change} (previous figure) for any causal reading.}

\projfig{0.70}{C_saving_by_quintile}{\textbf{Leg C, structure.} Median saving
rate by income quintile (Eurostat ICW, latest snapshot). The bottom quintile
\emph{dissaves}; saving climbs steeply to the top. Any rise in the aggregate rate
is mechanically a story about the upper part of the distribution.}

\projfig{0.70}{C2_expectations_by_income}{\textbf{Leg C, fear.} Expected
unemployment 12 months ahead by income group (ECB CES). Lower-income households
fear job loss most, yet (previous figure) can least afford to save---the equity
twist of the precautionary story.}

\projfig{0.74}{D_icw_slope}{\textbf{Leg C, over time.} The saving rate by
quintile across the few ICW reference years. Read as a structural shift between
snapshots, not as an annual path.}

% ============================================================================
\section{Econometric methodology}\label{sec:econ}
% ============================================================================

This section is deliberately detailed. The aim is to test leg~(A)
---does uncertainty \emph{drive} saving?---while taking seriously the statistical
properties of the series. All tests run on the aligned quarterly panel
$\{\text{saving}_t,\ \text{proxy}_t,\ \text{rate}_t,\ \text{inflation}_t\}$,
$t=$ 2000Q4$\ldots$2025Q4.

% ----------------------------------------------------------------------------
\subsection{Design principle: respect the integration order}
% ----------------------------------------------------------------------------
An earlier version of this analysis differenced \emph{every} series and fitted a
VECM. That was a mistake worth recording, because it motivates the current
design:
\begin{itemize}
  \item If a series is already stationary (I(0)), differencing it
  \emph{over-differences}: it discards the level information that carries the
  precautionary signal and biases impulse responses toward zero (a
  unit-root/non-invertible MA component is introduced).
  \item A VECM (and the Johansen test behind it) is only meaningful when
  \emph{all} series are I(1) and share cointegrating relationships. Feeding it
  stationary variables produces a ``long-run'' reading that is an artefact.
\end{itemize}
The current strategy therefore (1) classifies each series' order, (2) tests the
long-run relationship with a method valid for a \emph{mix} of I(0) and I(1)
regressors (ARDL bounds), and (3) runs Granger and the VAR/IRF in the
representation each series actually needs---levels for the stationary core. ARMA
is intentionally not used: it models a single series' own memory, not a
driver relationship.

% ----------------------------------------------------------------------------
\subsection{Step 1 --- Integration order from ADF and KPSS together}
% ----------------------------------------------------------------------------
We classify each series with two complementary unit-root tests, because they
have \emph{opposite} null hypotheses and so cross-check each other.

\paragraph{Augmented Dickey--Fuller (ADF).} Regress
\begin{equation}
  \Delta y_t = \alpha + \gamma\,y_{t-1}
             + \sum_{i=1}^{p}\delta_i\,\Delta y_{t-i} + \varepsilon_t,
  \qquad H_0:\ \gamma = 0\ \text{(unit root, non-stationary)} .
\end{equation}
Lag length $p$ is chosen by AIC. A \emph{small} $p$-value rejects the unit root
$\Rightarrow$ the series is stationary.

\paragraph{KPSS.} Decompose $y_t = \xi t + r_t + \varepsilon_t$ with a random
walk $r_t = r_{t-1}+u_t$ and test
\begin{equation}
  H_0:\ \sigma^2_u = 0\ \text{(level/trend stationary)} .
\end{equation}
Here the null is the \emph{opposite} of ADF's. We use the level-stationary form
(\code{regression="c"}). A \emph{large} $p$-value fails to reject stationarity.

\paragraph{Decision rule.} The series is confidently I(0) when ADF rejects
\emph{and} KPSS does not; confidently I(1) when ADF does not reject \emph{and}
KPSS does. When the two disagree (``\textsc{conflict}'') the order is genuinely
ambiguous---a hallmark of borderline or structurally-broken series---which is
exactly the case the long-run test below is chosen to tolerate. The collector
codes the order as I(0) iff ADF rejects at 5\%, and flags agreement with KPSS.

\begin{table}[htbp]\centering
\caption{Integration order, full sample (proxy $=$ GPR). $p$-values; ADF $H_0$:
unit root, KPSS $H_0$: stationarity.}
\label{tab:integ}
\renewcommand{\arraystretch}{1.15}
\begin{tabular}{@{}lccccc@{}}
\toprule
series & ADF $p$ (level) & KPSS $p$ (level) & ADF $p$ ($\Delta$) & order & agree? \\
\midrule
saving    & 0.017 & 0.087 & 0.000 & I(0) & yes \\
gpr       & 0.000 & 0.100 & 0.000 & I(0) & yes \\
rate      & 0.061 & 0.083 & 0.001 & I(1) & \textsc{conflict} \\
inflation & 0.652 & 0.100 & 0.000 & I(1) & \textsc{conflict} \\
\bottomrule
\end{tabular}
\end{table}

\noindent\textbf{Reading.} The \emph{core} variables---saving and the GPR
proxy---are clearly I(0): both tests agree. The two controls are ambiguous
(borderline I(1)). A system mixing I(0) and ambiguous-I(1) series is precisely
where blanket differencing or Johansen would mislead, and where ARDL bounds is
the right tool.

% ----------------------------------------------------------------------------
\subsection{Step 2 --- Long-run relationship: the ARDL bounds test}
% ----------------------------------------------------------------------------
We test for a \emph{level} (long-run) relationship between saving and the drivers
with the Pesaran--Shin--Smith (2001) bounds test, estimated through the
unrestricted (conditional) error-correction model. For dependent variable
$y_t=\text{saving}_t$ and regressors $x_t=(\text{proxy},\text{rate},
\text{inflation})$,
\begin{equation}
  \Delta y_t = c
    + \sum_{i=1}^{p-1}\phi_i \Delta y_{t-i}
    + \sum_{j}\sum_{l=0}^{q_j}\theta_{j l}\,\Delta x_{j,t-l}
    + \underbrace{\rho\,y_{t-1} + \sum_{j}\pi_j\,x_{j,t-1}}_{\text{lagged levels}}
    + \varepsilon_t .
  \label{eq:uecm}
\end{equation}
The bounds $F$-test has
\begin{equation}
  H_0:\ \rho = \pi_1 = \dots = \pi_k = 0
  \quad\text{(no long-run level relationship).}
\end{equation}

\paragraph{Why this test.} Its decisive advantage is that the critical values
come as a \emph{pair of bounds}: a lower bound computed as if all regressors were
I(0) and an upper bound as if all were I(1). The decision is therefore
\emph{order-agnostic}---valid for any mix of I(0) and I(1) regressors (it
requires only that no series is I(2)):
\begin{itemize}
  \item $F >$ upper bound $\Rightarrow$ reject $H_0$: a long-run relationship
  exists regardless of integration order;
  \item $F <$ lower bound $\Rightarrow$ cannot reject $H_0$: no long-run
  relationship;
  \item lower $\le F \le$ upper $\Rightarrow$ inconclusive.
\end{itemize}
Given the \textsc{conflict} flags in Table~\ref{tab:integ}, this is the correct
long-run test; Johansen is reported only for reference (below) and should be
ignored unless every series is I(1).

\paragraph{Implementation.} The forced-saving COVID quarters (2020) are dropped
as outliers. The model order is selected by minimising AIC over endogenous lags
$L\in\{1,\dots,4\}$ and exogenous order $O\in\{1,\dots,4\}$ (every exogenous
order $\ge 1$, as the UECM requires), with an unrestricted constant and no trend
(Pesaran et al.\ ``case III'').

\paragraph{Result (GPR).} The AIC-selected specification is $L=3$, $O=1$. The
bounds statistic is
\[
  F = 0.94 \quad\text{vs.\ the 5\% I(1) upper bound } 4.01,
\]
so $F$ lies \emph{below even the lower bound}: we cannot reject $H_0$. There is
\textbf{no statistical evidence of a long-run level relationship} between the
saving rate and the drivers over this sample. (For reference only, the Johansen
trace implies rank 2, but with mixed integration orders this is uninformative and
disregarded.)

\begin{table}[htbp]\centering
\caption{ARDL bounds critical values (Pesaran--Shin--Smith, case III). The
test statistic $F=0.94$ is below the lower bound at every level.}
\renewcommand{\arraystretch}{1.15}
\begin{tabular}{@{}lcc@{}}
\toprule
significance & lower bound, I(0) & upper bound, I(1) \\
\midrule
10\% & 2.46 & 3.52 \\
 5\% & 2.88 & 4.01 \\
 1\% & 3.78 & 5.05 \\
\bottomrule
\end{tabular}
\end{table}

% ----------------------------------------------------------------------------
\subsection{Step 3 --- Granger causality}
% ----------------------------------------------------------------------------
Granger causality asks whether past values of the proxy improve the prediction of
saving \emph{beyond} saving's own past. We test it in the stationary
representation of each series (levels for the I(0) core, here), drop COVID 2020,
and scan lags $1$--$4$:
\begin{equation}
  \text{saving}_t = \alpha + \sum_{i=1}^{m}\beta_i\,\text{saving}_{t-i}
    + \sum_{i=1}^{m}\eta_i\,\text{proxy}_{t-i} + u_t,
  \qquad H_0:\ \eta_1=\dots=\eta_m=0 .
\end{equation}
An $F$-test $p<0.05$ on the $\{\eta_i\}$ means the proxy Granger-causes saving.

\paragraph{Result (GPR, full sample).} The $p$-values for ``uncertainty $\to$
saving'' are $0.91,\,0.44,\,0.71,\,0.72$ at lags $1$--$4$: \textbf{no Granger
causality} from uncertainty to saving at any lag. (The reverse direction,
saving $\to$ uncertainty, is significant at lags 3--4, but a household saving
rate ``predicting'' a global geopolitical-risk index is not economically
meaningful and most likely reflects common low-frequency movement.) Two cautions
apply: Granger causality is \emph{predictive}, not structural; and on $\sim$100
quarters with a large COVID hole the test has limited power, so a null is weak
evidence of absence, not evidence of a null.

% ----------------------------------------------------------------------------
\subsection{Step 4 --- Structural VAR and impulse responses}
% ----------------------------------------------------------------------------
To trace the dynamic response of saving to an uncertainty shock while
controlling for the confounders, we estimate a VAR and compute orthogonalised
impulse-response functions (IRFs).

\paragraph{Specification.} With $Y_t=(\text{proxy},\text{rate},\text{inflation},
\text{saving})'$,
\begin{equation}
  Y_t = c + \sum_{i=1}^{p} A_i\,Y_{t-i} + B\,D_t + u_t,
  \qquad u_t \sim (0,\Sigma),
\end{equation}
where $D_t$ are dummies for the forced-saving COVID quarters (2020Q2--2021Q1) so
that the pandemic spike does not contaminate the dynamics. Because the core
series are I(0) (Table~\ref{tab:integ}), the VAR is estimated \emph{in levels},
and the IRF is read directly as a level response---no cumulation, no
over-differencing. (If the core were not both I(0), the code estimates in first
differences and cumulates the IRF back to levels.) Lag length $p$ is chosen by
AIC and capped at $4$ to keep the responses stable on a short sample.

\paragraph{Identification.} Shocks are orthogonalised by a Cholesky
factorisation $\Sigma = PP'$ with the recursive ordering
\[
  \text{proxy} \;\to\; \text{rate} \;\to\; \text{inflation} \;\to\; \text{saving}.
\]
The identifying assumption is that a variable ordered earlier may affect those
ordered later \emph{within the quarter}, but not the reverse. Placing uncertainty
\emph{first} treats it as the most exogenous block---a geopolitical-risk index is
not plausibly driven by euro-area household saving within a quarter---and placing
saving \emph{last} lets it react contemporaneously to all three drivers. This is
the natural ordering for the precautionary question; like any recursive scheme it
is an assumption, not a fact, and is the main identification caveat.

\paragraph{Result (GPR).} The response of saving to a $+1$ standard-deviation GPR
shock is \emph{positive but statistically insignificant}:
\begin{center}
\begin{tabular}{@{}lcc@{}}
\toprule
sample & peak response & 95\% CI at peak \\
\midrule
full (2000Q4--2025Q4) & $+0.108$ pp at 0q & $[-0.015,\ +0.230]$ \\
subsample (2010+)     & $+0.161$ pp at 6q & $[-0.019,\ +0.340]$ \\
\bottomrule
\end{tabular}
\end{center}
The sign is \emph{consistent with} precaution---more uncertainty, more
saving---but the confidence interval includes zero in both samples, so the effect
is not statistically distinguishable from no effect. Figure~\ref{fig:A2_irf_saving_to_uncertainty}
plots the full-sample IRF with its band.

\projfig{0.66}{A2_irf_saving_to_uncertainty}{Orthogonalised impulse response of
the saving rate to a $+1$ s.d.\ GPR shock (levels VAR; controls: ECB rate and
HICP inflation; COVID dummies). The 95\% band straddles zero throughout.}

% ----------------------------------------------------------------------------
\subsection{Step 5 --- Robustness: subsample and proxy choice}
% ----------------------------------------------------------------------------
Two robustness dimensions are built in. First, the key tests are re-run on a
\textbf{post-2010 subsample} (64 quarters)---the period most relevant to the
modern saving puzzle---with the same qualitative result (positive, insignificant
IRF; no Granger causality). Second, the entire analysis can be re-run with the
\textbf{EPU} proxy instead of GPR (\code{--proxy epu}); EPU yields the same
narrative (a positive but insignificant response), confirming the conclusion is
not an artefact of one proxy.

% ----------------------------------------------------------------------------
\subsection{Step 6 --- Proxy and timing: lead/lag cross-correlation}
\label{sec:timing}
% ----------------------------------------------------------------------------
Finally we ask, descriptively, \emph{which} proxy and at \emph{what} lead/lag
best tracks saving. After $z$-scoring both series we compute the
cross-correlation at quarterly shifts $k=-8,\dots,+8$,
\begin{equation}
  \rho_k = \operatorname{corr}\!\big(z(\text{saving})_t,\;
                                     z(\text{proxy})_{t-k}\big),
\end{equation}
where $k>0$ means the proxy \emph{leads} saving. We report it on both levels and
first differences ($\Delta$), because the levels of trending series mechanically
inflate correlations, and we drop COVID 2020.

\paragraph{Result.} In levels, GPR attains the larger peak correlation
($|\rho|=0.42$) than EPU ($0.35$), which is why GPR is the preferred headline
proxy. The peak occurs at $k=-4$ for GPR (i.e.\ the proxy is contemporaneous-to-
slightly-\emph{lagging} rather than cleanly leading), and in differences the
peaks are weaker ($-0.32$ for GPR at $k=-2$). The lead/lag evidence is therefore
suggestive of co-movement but does \emph{not} pin down a clean
``uncertainty-leads-saving'' timing---consistent with the insignificant formal
tests. Figure~\ref{fig:A3_proxy_leadlag} shows the full profiles.

\projfig{0.92}{A3_proxy_leadlag}{Lead/lag cross-correlation of the saving rate
with each proxy, in levels (left) and changes (right). $k>0$: proxy leads saving.
COVID 2020 dropped.}

% ============================================================================
\section{Did rates or fear drive it? The follow-the-money test}\label{sec:followmoney}
% ============================================================================

The methodology of Section~\ref{sec:econ} targets the \emph{level} of the saving
rate and returns nulls: with so few quarters and three drivers moving together
after 2022, no test can cleanly assign the rise to precaution. This section asks a
different, sharper question --- about the \emph{composition} of saving rather than
its level --- where the precautionary and the ``yield-chasing'' (intertemporal-%
substitution) explanations make \emph{opposite} predictions and can therefore be
separated. The analysis lives in \code{extension\_follow\_money/}.

\subsection{The discriminating idea: where the money goes}
Fear and yield make opposite predictions about \emph{where} households park new
saving. Precaution favours liquid, instant-access assets (the buffer must be
reachable in an emergency); yield-chasing favours higher-paying but less-liquid
ones. We split households' annual net acquisition of financial assets (Eurostat
\dataset{nasa\_10\_f\_tr}, sector S14\_S15, euro area, current-price flows) into
two buckets and track the \emph{tilt} between them:
\begin{align*}
\text{instant-access} &= \text{currency (F21)} + \text{overnight deposits (F22)},\\
\text{locked-for-yield} &= \text{time deposits (F29)} + \text{bonds (F3)},\\
\text{tilt}_t &= \text{locked-for-yield}_t - \text{instant-access}_t \quad(\text{EUR bn/yr}).
\end{align*}
A positive tilt means that year's saving leaned toward yield; negative, toward
instant cash.

\subsection{Linking the asset flows back to the saving rate}
The tilt is a \emph{composition} measure, so it does not map to the saving rate on
its own. The bridge is an accounting identity: when a household saves, the money
flows into financial assets, into housing, or into paying down debt; the part going
into financial assets is exactly the \emph{sum of all the asset-type flows we
decompose}, and the tilt is \emph{where within that sum} the money leans. So if the
decomposition is legitimate, combining all the buckets should track the saving
rate.

It does. Summing the asset flows and scaling by household gross disposable income
($B6G$, Eurostat \dataset{nasa\_10\_nf\_tr}) gives a series that correlates
$+0.87$ with the saving rate in levels and $+0.73$ as a share of income; both
spike in 2020, dip in 2022--23, and recover. As a reconciliation, rebuilding the
rate from the raw income accounts (gross saving $B8G$ over disposable income
$B6G$) reproduces the published saving rate almost exactly (correlation $+1.00$;
Table~\ref{tab:reconcile}). The combined-flows line sits a few points \emph{below}
the saving rate --- the wedge is mainly housing investment, which saving funds but
is not a financial asset --- yet the two clearly co-move
(Figure~\ref{fig:total_vs_saving}). This places the composition analysis firmly
inside the saving story: the tilt describes \emph{how the saving in the headline
rate was deployed}.

\begin{table}[htbp]\centering
\caption{Reconciliation: the rate rebuilt from the income accounts ($B8G/B6G$)
matches the published saving rate, and the combined financial-asset flow tracks it
(all in \% of disposable income).}
\label{tab:reconcile}
\renewcommand{\arraystretch}{1.15}
\begin{tabular}{@{}lccc@{}}
\toprule
year & combined flows & saving rate ($B8G/B6G$) & saving rate (published) \\
\midrule
2020 & 15.0 & 19.7 & 19.5 \\
2022 &  8.8 & 13.6 & 13.5 \\
2023 &  6.9 & 14.4 & 14.2 \\
2025 &  9.3 & 15.0 & 14.8 \\
\bottomrule
\end{tabular}
\end{table}

\extfig{0.9}{total_vs_saving}{Combining all the asset flows reproduces the
saving-rate path. Navy: the published household saving rate. Green dashed: the sum
of all net financial-asset flows, as a share of disposable income. They co-move
($\mathrm{corr}=+0.73$); the roughly constant gap is housing investment.}

\subsection{Descriptive evidence}
The destination of saving flipped after the ECB began raising rates in mid-2022.
The share of yearly financial saving going into instant-access cash fell from
about 66\% (2015--19) to about 2\% (2022--25), while net household \emph{bond}
purchases went from negative (net selling of roughly EUR 40--50 bn/yr in 2019--21)
to +EUR 96 bn (2022) and +EUR 311 bn (2023). Households went from shunning bonds
to buying them in record size precisely as yields rose --- the reach-for-yield
pattern, and the opposite of the 2020 COVID episode, when saving piled into
instant-access cash (the precautionary pattern). The tilt correlates $+0.79$ with
the ECB policy rate (annual). Figures~\ref{fig:tilt_vs_saving}
and~\ref{fig:tilt_vs_rates} show the picture.

\extfig{0.95}{tilt_vs_saving}{Composition tilt (bars; locked-for-yield minus
instant-access, EUR bn/yr) against the household saving rate (line). The saving
\emph{level} stays high throughout --- peaking in the COVID year, when the tilt is
most negative (cash) --- while the \emph{composition} flips to yield (red) once
rates rise. Regimes shaded as in chart~A: the zero-lower-bound era, COVID forced
saving, and the post-2022 period of interest.}

\extfig{0.95}{tilt_vs_rates}{The composition tilt (bars) against the ECB policy
rate (line); the two move together, $\mathrm{corr}=+0.79$ over 2002--2025.}

\subsection{A model comparison: rates/inflation versus precaution}
Because the tilt is signed, the two stories give opposite, testable predictions.
We regress the tilt on the ECB rate, inflation, and an uncertainty proxy (the
Geopolitical Risk index, GPR), each standardised to a $z$-score so the
coefficients are comparable (the effect of a one-standard-deviation move, in
EUR bn):
\begin{equation}
\text{tilt}_t = \beta_0 + \beta_1 \widetilde{r}_t + \beta_2 \widetilde{\pi}_t
              + \beta_3 \widetilde{u}_t + \varepsilon_t,
\label{eq:tilt}
\end{equation}
with $\widetilde r$ the rate, $\widetilde\pi$ inflation, $\widetilde u$
uncertainty. The competing hypotheses are sharp:
\[
\text{yield-chasing} \;\Rightarrow\; \beta_1 > 0, \qquad
\text{precaution} \;\Rightarrow\; \beta_3 < 0
\]
(more uncertainty should push saving \emph{toward} instant cash, a more negative
tilt). Standard errors are heteroskedasticity-robust (HC3), appropriate for the
small sample.

\paragraph{Letting the drivers compete.} Rather than lean on information criteria,
we simply let the three drivers compete in a single regression and read the result
two ways: (i) which coefficient is statistically significant once all are
included, and (ii) how much of the year-to-year variation each story explains on
its own ($R^2$). The regression table (Table~\ref{tab:compecon}) does the rest.

\subsection{Results}
Table~\ref{tab:compecon} is the horse race for the tilt (annual, $n=24$). The
reading is unambiguous. With all drivers included, \emph{only the ECB rate is
significant} ($+304$, $p\approx0.02$); inflation is not, and uncertainty is not
--- and uncertainty carries the \emph{wrong} sign for precaution. The lopsided fit
says the same: the rates+inflation model explains about $74\%$ of the tilt's
year-to-year variation (adjusted $R^2$), the precaution model about $2\%$. Using
net bond purchases as the dependent variable gives the same verdict (rate $+62$,
$p\approx0.03$); there uncertainty even looks significant \emph{on its own}
($+43$, $p\approx0.02$) but collapses to zero once the rate is included --- a
textbook confounder, with uncertainty merely standing in for rates.

\begin{table}[htbp]\centering
\caption{The composition tilt, explained: a horse race (annual, 2002--2025).
Dependent variable: net flow into locked-for-yield minus instant-access
(EUR bn/yr). Drivers standardised to $z$-scores, so each coefficient is EUR bn per
one-standard-deviation move. HC3 robust standard errors in parentheses.}
\label{tab:compecon}
\renewcommand{\arraystretch}{1.05}
\begin{tabular}{@{}lccc@{}}
\toprule
 & \multicolumn{3}{c}{\emph{Dependent variable:} composition tilt}\\
\cmidrule(lr){2-4}
 & (1) rates$+$infl & (2) precaution & (3) all\\
\midrule
ECB rate          & $278.3^{**}$ &            & $303.6^{**}$\\
                  & $(123.3)$    &            & $(125.6)$\\
Inflation         & $160.9$      &            & $193.5$\\
                  & $(231.0)$    &            & $(194.9)$\\
Uncertainty (GPR) &              & $103.3$    & $-95.5$\\
                  &              & $(65.6)$   & $(83.9)$\\
Constant          & $-123.5^{*}$ & $-123.5$   & $-123.5^{*}$\\
                  & $(70.9)$     & $(88.6)$   & $(66.7)$\\
\midrule
$N$               & 24    & 24   & 24\\
$R^2$             & 0.76  & 0.06 & 0.80\\
Adjusted $R^2$    & 0.74  & 0.02 & 0.77\\
\bottomrule
\multicolumn{4}{@{}l}{\footnotesize $^{*}\,p<0.1$,\; $^{**}\,p<0.05$,\; $^{***}\,p<0.01$.}\\
\end{tabular}
\end{table}

\subsection{Robustness and scope}
Regressing levels of trending series can manufacture significance, so we
re-estimate \eqref{eq:tilt} in first differences. The rate coefficient keeps the
right (positive) sign but loses significance ($p\approx0.13$) on $n=23$ noisy
annual changes; the \emph{ranking} is unchanged, though --- the rate still
dominates the fit while uncertainty stays insignificant and wrong-signed. Two
caveats bound the claim. The
sample is small ($\sim$24 annual euro-area observations) and the 2023 bond surge
is a high-leverage point; robust errors and the differenced re-run mitigate but do
not remove this. And, most important, this test concerns the \emph{composition} of
saving --- \emph{how} households deployed it --- not the \emph{level} of the saving
rate: it shows the reallocation of saving is far more likely a reaction to rates
and inflation than to precaution, but does not by itself explain why the aggregate
rate rose, which stays overdetermined (Section~\ref{sec:econ}).

% ============================================================================
\section{Interpretation}\label{sec:interpretation}
% ============================================================================

Putting the legs together:
\begin{itemize}
  \item \textbf{Leg A (time series).} Descriptively strong: saving is elevated
  and rose alongside uncertainty after 2022 (Fig.~\ref{fig:A_saving_vs_uncertainty}).
  But once the ECB rate and inflation are controlled for, the formal tests
  cannot establish a precautionary channel: no long-run relationship (ARDL
  $F=0.94$), no Granger causality, and a positive-but-insignificant impulse
  response. The honest conclusion is \emph{co-movement consistent with, but not
  proof of, precaution.}
  \item \textbf{Leg B (cross-country).} Suggestive and confounded: the slope of
  saving-change on energy-shock is positive but noisy, and the level map mostly
  reflects long-standing North--South habits.
  \item \textbf{Leg C (distribution).} The clearest structural fact: the bottom
  quintile dissaves and the top saves, so any aggregate rise is concentrated
  among households with slack---while lower-income households report the most
  fear. The same uncertainty produces a bigger buffer at the top and forgone
  consumption at the bottom.
\end{itemize}
The overall verdict is intentionally measured: the precautionary story is
\emph{plausible and consistent with the data}, but the time-series evidence is
not strong enough to claim a clean causal channel---which is why the descriptive
charts and the econometrics are presented together rather than one standing in
for the other. On the narrower question of \emph{how} households saved, however,
the evidence is sharper: the follow-the-money test (Section~\ref{sec:followmoney})
finds the post-2022 reallocation toward yield is far more likely a reaction to
rates and inflation than to precaution --- a statement the level data alone cannot
support.

% ============================================================================
\section{Limitations and caveats}\label{sec:limits}
% ============================================================================
\begin{itemize}
  \item \textbf{Power.} The aligned sample is short ($\sim$101 quarters) and
  contains a large COVID structural break; time-series tests on such samples have
  low power, so null results are weak evidence of absence.
  \item \textbf{Identification.} The IRF rests on a recursive (Cholesky)
  ordering; a different ordering or a sign-restriction scheme could shift the
  point estimates (though the wide bands suggest not the significance).
  \item \textbf{Proxies.} Uncertainty is latent; GPR and EPU are imperfect,
  newspaper-based measures.
  \item \textbf{Gross vs.\ net saving.} The headline target is the \emph{gross}
  household saving rate; cross-source comparisons should note this.
  \item \textbf{Distributional timeliness.} ICW is experimental and only
  $\sim$5-yearly; it shows structure, not the post-2022 dynamics. The OECD annual
  panel that would show year-by-year movement is not currently retrievable from
  the queried flow.
  \item \textbf{Editorial exposure grouping.} The red/blue Russia-energy split in
  leg~B is a simplification; the continuous energy-inflation axis is the rigorous
  version.
  \item \textbf{Composition vs.\ level.} The follow-the-money test
  (Section~\ref{sec:followmoney}) speaks to \emph{how} saving was allocated, not
  \emph{why} the aggregate rate rose, and rests on $\sim$24 annual observations
  with a high-leverage 2023.
  \item \textbf{Vintages.} All series are pulled live and revised by providers;
  regenerate before citing any number.
\end{itemize}

% ============================================================================
\appendix
\section{Reproduction details}
% ============================================================================

\subsection{File tree}
\begin{lstlisting}
01_collect_data.py     # STEP 1  data collection  -> data/*.csv
02_make_figures.py     # STEP 2  figures          -> figures/*.png
03_econometrics.py     # STEP 3  econometrics      -> results.md, A2, A3
run_all.sh             # pip install + run 01->02->03 (classic Python, no venv)
requirements.txt       # pinned dependencies
data/                  # generated tidy CSVs
figures/               # generated charts
extensions/            # identification extensions (panel FE, event study, ...)
extension_follow_money/# follow-the-money composition tests + econometrics (Sec. 6)
findings.md            # narrative write-up of the results
documentation.tex      # this document
legacy/                # superseded original scripts (safe to delete)
\end{lstlisting}

\subsection{Commands}
\begin{lstlisting}
# one-shot, fully reproducible
bash run_all.sh

# or stage by stage (after: pip install -r requirements.txt)
python 01_collect_data.py
python 02_make_figures.py
python 03_econometrics.py            # GPR (default)
python 03_econometrics.py --proxy epu --substart 2010
\end{lstlisting}

\subsection{Data dictionary (generated CSVs)}
\begin{center}
\renewcommand{\arraystretch}{1.2}
\begin{tabular}{@{}p{5.3cm} p{3.0cm} p{6.0cm}@{}}
\toprule
file & frequency & contents \\
\midrule
\code{ea\_saving\_rate\_quarterly.csv} & quarterly & date, saving (\% GDI) \\
\code{gpr.csv} & monthly & date, GPR index \\
\code{fred\_eu\_epu.csv} & monthly & date, EU EPU index \\
\code{ecb\_rate.csv} & daily & date, ECB deposit rate (\%) \\
\code{ea\_inflation.csv} & monthly & date, HICP YoY (\%) \\
\code{country\_saving\_annual.csv} & annual & geo $\times$ year saving pivot \\
\code{country\_energy\_peak\_2022.csv} & 2022 scalar & geo, peak energy infl.\ (\%) \\
\code{scatter\_energy\_vs\_saving.csv} & cross-section & geo, saving rise, energy peak \\
\code{saving\_rate\_by\_quintile.csv} & snapshot & q, value, geo, year \\
\code{icw\_quintile\_panel.csv} & snapshots & year, q, value, geo \\
\code{ces\_unemp\_expectations.csv} & latest & group, expected unemployment (\%) \\
\code{econometrics\_results.md} & --- & full text log of step 3 \\
\bottomrule
\end{tabular}
\end{center}

\vfill
\noindent\rule{\linewidth}{0.4pt}\\[2pt]
{\small All numerical results in this document come from a clean end-to-end run
and are reproducible via \code{run\_all.sh}. Because the underlying series are
revised by their providers, regenerate the data and re-compile before citing.}

\end{document}
